\section{Analysis and Negative Results}
\label{sec:analysis}

The negative results are part of the contribution; they justify the bounded scope.

\begin{itemize}
  \item[(i)] \textbf{The original graph backbone fails the task.} A GraphCMINet-style network is near-chance
  under strict source-only training and is not repairable by small training tweaks, so its (strong,
  controllable) leakage cannot support a task-vs-leakage tradeoff.
  \item[(ii)] \textbf{The protocol is learnable.} Known-good convolutional decoders
  (EEGNet/ShallowConvNet/DeepConvNet) and a DGCNN all clear a source-accuracy floor on the same fold,
  showing the bottleneck was the specific graph backbone, not the data/protocol.
  \item[(iii)] \textbf{Dynamic-edge designs overfit.} The tested dynamic-adjacency backbones overfit
  strongly (train $\approx 1.0$, source $\approx$ chance) regardless of temporal stem, and carry leakage;
  this is consistent with a per-sample adjacency $A(x)$ acting as a subject-fingerprint channel, but these
  experiments do \textbf{not} causally isolate $A(x)$ as the only source of memorization. Either way,
  \textbf{dynamic-edge CIGL is unsupported here and edge-CMI is out of scope}.
\end{itemize}

The static-adjacency DGCNN is the only graph-compatible backbone that both learns the task and exposes
usable graph/node objects, which is why CIGL is framed as graph/node only on this backbone. Table~5
summarizes the negative results, and Figure~F4 shows the decision flow that fixed the scope.
